% Lecture 4 notes: \input by notes.tex, \subimport'ed by the semester book.
% Exam questions (экзамен): 9 (метод исчисления индексов, оценка сложности).
\begin{lecture}{4}{Метод исчисления индексов}

\section{Зачем выходить за рамки универсальных алгоритмов}

В лекциях 1--3 мы решали задачу дискретного логарифмирования (DLP) $g^x = h$ алгоритмами,
которые работают в \emph{любой} циклической группе: перебором, алгоритмом Шенкса,
$\rho$-методом Полларда, а для группы составного порядка "--- алгоритмом Полига"--~Хеллмана. Все
они \emph{универсальны}: используют только групповые операции и сравнение элементов. По теореме
Шупа (лекция~1, \cite{shoup-generic}) любой универсальный алгоритм выполняет
$\Omega(\sqrt q)$ групповых операций, где $q$ "--- наибольший простой делитель порядка группы.
Если бы для $\Znx{p}$ не было ничего лучше, для 128-битной стойкости протокола Диффи"--~Хеллмана
хватило бы простого $p$ длиной около $256$ бит с большим простым делителем $p - 1$.

На практике же рекомендуемая длина $p$ "--- не меньше $2048$--$3072$ бит. Причина в том, что
элементы $\Znx{p}$ "--- это \emph{числа}, и их можно раскладывать на простые множители.
Дискретный логарифм "--- гомоморфизм (лекция~1):
\[
  \log_g(ab) \equiv \log_g a + \log_g b \pmod{p - 1},
\]
поэтому, если $h \equiv 2^{3} \cdot 5^{2} \pmod p$, то $\log_g h \equiv 3\log_g 2 + 2\log_g 5$, и
задача сводится к логарифмам \emph{малых простых}. Логарифмы малых простых, в свою очередь,
находятся из линейной системы, которую дают разложения степеней $g^k \bmod p$. Эта идея восходит
к М.~Крайчику (1922); как алгоритм с доказанной субэкспоненциальной оценкой её оформил
Л.~Адлеман~\cite{adleman-index-calculus} (1979). Метод называют \emph{методом исчисления
индексов} (index calculus): «индекс» "--- старое название дискретного логарифма,
$\ind_g h = \log_g h$.

Центральные понятия лекции "--- \emph{гладкие числа} (числа без больших простых делителей) и
\emph{субэкспоненциальная} сложность, промежуточная между полиномиальной и экспоненциальной.
Те же идеи лежат в основе алгоритмов факторизации (квадратичное решето, лекция~6).

\section{Гладкие числа}

\begin{definition}
Пусть $B \ge 2$. Натуральное число $m$ называется \emph{$B$-гладким}, если все его простые
делители не превосходят $B$ (число $1$ считается $B$-гладким). Через $\psi(x, B)$ обозначается
количество $B$-гладких чисел $m \le x$.
\end{definition}

Например, $5$-гладкие числа до $30$: $1, 2, 3, 4, 5, 6, 8, 9, 10, 12, 15, 16, 18, 20, 24, 25, 27,
30$, т.\,е. $\psi(30, 5) = 18$; далее $\psi(100, 5) = 34$, $\psi(700, 7) = 121$. Проверка
$B$-гладкости числа $m$ \emph{пробным делением} на простые $\ell \le B$ требует $\pi(B)$ делений
(плюс деления, отщепляющие кратные множители, "--- их не больше $\log_2 m$).

Насколько часто встречаются гладкие числа? Удобно параметризовать границу как $B = x^{1/u}$, т.\,е.
\[
  u = \frac{\ln x}{\ln B}
\]
"--- во сколько раз число $x$ «длиннее» границы $B$. Начнём с простой нижней оценки.

\begin{proposition}\label{l4:psi-lower}
Пусть $k \ge 1$ "--- целое, $B^k \le x$. Тогда
\[
  \psi(x, B) \ \ge\ \binom{\pi(B) + k - 1}{k} \ \ge\ \frac{\pi(B)^k}{k!}.
\]
В частности, если $u = \ln x / \ln B$ "--- целое и $B \ge 17$, то
$\psi(x, B)/x \ge (u \ln B)^{-u} = (\ln x)^{-u}$.
\end{proposition}

\begin{proof}
Рассмотрим произведения $\ell_1 \ell_2 \cdots \ell_k$ простых $\ell_1 \le \ell_2 \le \cdots \le
\ell_k \le B$. Каждое из них $B$-гладкое и не превосходит $B^k \le x$, а по основной теореме
арифметики различные неубывающие наборы дают различные числа. Число таких наборов "--- число
сочетаний с повторениями из $\pi(B)$ по $k$, т.\,е. $\binom{\pi(B) + k - 1}{k}$. Далее,
\[
  \binom{\pi(B) + k - 1}{k} = \frac{\pi(B)(\pi(B) + 1) \cdots (\pi(B) + k - 1)}{k!}
  \ge \frac{\pi(B)^k}{k!}.
\]
Для второго утверждения возьмём $k = u$, так что $B^u = x$. Известно, что $\pi(B) > B/\ln B$ при
$B \ge 17$ (Россер"--~Шёнфельд; асимптотически это следует из теоремы о распределении простых
$\pi(B) \sim B/\ln B$). Так как $u! \le u^u$,
\[
  \frac{\psi(x, B)}{x} \ge \frac{\pi(B)^u}{u!\,B^u} \ge \frac{1}{u^u (\ln B)^u}
  = (u \ln B)^{-u} = (\ln x)^{-u}. \qedhere
\]
\end{proof}

Оценка груба: она учитывает только числа ровно из $u$ простых множителей. Правильный порядок
даёт следующая глубокая теорема.

\begin{theorem}[Кэнфилд"--~Эрдёш"--~Померанс~\cite{canfield-erdos-pomerance}, без доказательства]
\label{l4:cep}
Для любого $\varepsilon > 0$ равномерно по $3 \le u \le (1 - \varepsilon) \ln x / \ln\ln x$
\[
  \psi\bigl(x, x^{1/u}\bigr) = x \cdot u^{-u(1 + o(1))}, \qquad x \to \infty.
\]
\end{theorem}

Иначе говоря, \emph{доля $B$-гладких чисел среди чисел до $x$ примерно равна $u^{-u}$}.
Для фиксированного $u$ известна и точная асимптотика $\psi(x, x^{1/u}) \sim \rho(u)\, x$, где
$\rho$ "--- функция Дикмана ($\rho(u) = 1$ при $u \le 1$, $\rho(u) = 1 - \ln u$ при $1 \le u \le
2$, $\rho(u) = u^{-u(1 + o(1))}$). Множитель $1 + o(1)$ в показателе существен: при $x = 10^6$,
$B = 100$ ($u = 3$) точное значение $\psi(10^6, 100) = 72\,271$, тогда как $\rho(3)\, x \approx
48\,600$, а $u^{-u} x \approx 37\,000$. Для оценок сложности, однако, важен лишь порядок
$u^{-u}$, а он теоремой~\ref{l4:cep} описан правильно.

\checkyourself{Найдите $\psi(50, 3)$. Сравните с нижней оценкой
утверждения~\ref{l4:psi-lower} при $k = 3$. Почему в утверждении~\ref{l4:psi-lower} нужно
требовать неубывание $\ell_1 \le \cdots \le \ell_k$?}

\section{\texorpdfstring{$L$}{L}-нотация и субэкспоненциальная сложность}

Длина входа задачи DLP в $\Znx{p}$ "--- это $\ln p$ (с точностью до множителя). Алгоритм
\emph{полиномиален}, если его сложность $(\ln p)^{\bigO(1)}$, и \emph{экспоненциален}, если она
равна $p^{c}$ (как $\sqrt p$ у алгоритма Шенкса). Промежуточные функции удобно записывать так.

\begin{definition}
Для $0 \le \alpha \le 1$ и $c > 0$
\[
  L_x[\alpha, c] = \exp\Bigl( (c + o(1)) (\ln x)^{\alpha} (\ln\ln x)^{1 - \alpha} \Bigr),
  \qquad x \to \infty.
\]
Часто пишут $L(x) = L_x[1/2, 1] = e^{\sqrt{\ln x \ln\ln x}}$ (без $o(1)$). Алгоритм сложности
$L_p[\alpha, c]$ с $0 < \alpha < 1$ называется \emph{субэкспоненциальным}.
\end{definition}

При $\alpha = 0$ получаем $L_x[0, c] = (\ln x)^{c + o(1)}$ "--- полиномиальную сложность, при
$\alpha = 1$ "--- $L_x[1, c] = x^{c + o(1)}$, экспоненциальную. Слагаемое $o(1)$ «поглощает»
полиномиальные множители:

\begin{lemma}\label{l4:Lprops}
Пусть $0 < \alpha < 1$. Тогда
\begin{enumerate}
  \item $L_x[\alpha, a] \cdot L_x[\alpha, b] = L_x[\alpha, a + b]$ и
    $L_x[\alpha, a] + L_x[\alpha, b] = L_x[\alpha, \max(a, b)]$;
  \item $(\ln x)^{k} \cdot L_x[\alpha, a] = L_x[\alpha, a]$ для любого фиксированного $k$;
  \item $L_x[\alpha, a] = o(x^{\varepsilon})$ для любого $\varepsilon > 0$ и
    $(\ln x)^{k} = o(L_x[\alpha, a])$ для любого $k$.
\end{enumerate}
\end{lemma}

\begin{proof}
Обозначим $\lambda = (\ln x)^{\alpha}(\ln\ln x)^{1 - \alpha}$; при $0 < \alpha < 1$
\[
  \frac{\lambda}{\ln\ln x} = \Bigl(\frac{\ln x}{\ln\ln x}\Bigr)^{\alpha} \to \infty, \qquad
  \frac{\lambda}{\ln x} = \Bigl(\frac{\ln\ln x}{\ln x}\Bigr)^{1 - \alpha} \to 0.
\]
(1) Показатели складываются: $(a + o(1))\lambda + (b + o(1))\lambda = (a + b + o(1))\lambda$.
Для суммы при $a \ge b$:
$L[\alpha, a] \le L[\alpha, a] + L[\alpha, b] \le 2 L[\alpha, a]$, и $\ln 2 = o(\lambda)$.
(2) $\ln\bigl((\ln x)^k\bigr) = k \ln\ln x = o(\lambda)$, т.\,е. множитель меняет показатель
на $o(1) \cdot \lambda$. (3) Отношение $L_x[\alpha, a] / x^{\varepsilon} =
\exp\bigl((a + o(1))\lambda - \varepsilon \ln x\bigr) \to 0$, так как $\lambda = o(\ln x)$;
аналогично $k \ln\ln x - (a + o(1))\lambda \to -\infty$.
\end{proof}

Свяжем $L$-нотацию с гладкостью. Ключевая выкладка всей лекции "--- следующая лемма.

\begin{lemma}\label{l4:uu}
Пусть $B = L_x[1/2, c]$ и $u = \ln x / \ln B$. Тогда $u^{u} = L_x[1/2, \tfrac{1}{2c}]$. Как
следствие (при условии теоремы~\ref{l4:cep}), доля $B$-гладких чисел среди чисел до $x$ равна
\[
  \frac{\psi(x, B)}{x} = L_x\Bigl[\frac12, -\frac{1}{2c}\Bigr]
  := \exp\Bigl(-\bigl(\tfrac{1}{2c} + o(1)\bigr)\sqrt{\ln x \ln\ln x}\Bigr).
\]
\end{lemma}

\begin{proof}
Положим $\Lambda = \ln x$. Тогда $\ln B = (c + o(1))\sqrt{\Lambda \ln \Lambda}$, и
\[
  u = \frac{\Lambda}{\ln B} = \frac{1 + o(1)}{c}\sqrt{\frac{\Lambda}{\ln \Lambda}}, \qquad
  \ln u = \tfrac12 \ln\Lambda - \tfrac12 \ln\ln\Lambda + \bigO(1) = \bigl(\tfrac12 + o(1)\bigr)\ln\Lambda.
\]
Перемножая,
\[
  \ln(u^u) = u \ln u = \frac{1 + o(1)}{c}\sqrt{\frac{\Lambda}{\ln\Lambda}} \cdot
  \bigl(\tfrac12 + o(1)\bigr) \ln\Lambda = \Bigl(\frac{1}{2c} + o(1)\Bigr)\sqrt{\Lambda \ln\Lambda},
\]
что и означает $u^u = L_x[1/2, 1/(2c)]$. Значение $u \asymp \sqrt{\Lambda / \ln\Lambda}$ лежит в
пределах теоремы~\ref{l4:cep} (там допускается $u$ до $(1 - \varepsilon)\Lambda / \ln\Lambda$), и
$\psi(x, B)/x = u^{-u(1 + o(1))} = L_x[1/2, -1/(2c)]$.
\end{proof}

Итак, если граница $B$ «субэкспоненциальна», то и доля гладких чисел убывает лишь
субэкспоненциально. Увеличивая $B$, мы чаще встречаем гладкие числа, но платим большим
количеством неизвестных; ниже мы найдём оптимальный баланс.

\checkyourself{Покажите, что при \emph{фиксированном} $B$ (например, $B = 100$) число
$B$-гладких чисел до $x$ растёт лишь полиномиально по $\ln x$, т.\,е. почти все числа не
$B$-гладкие. Указание: $\psi(x, B) \le (1 + \log_2 x)^{\pi(B)}$. Почему поэтому граница $B$ должна
расти вместе с $p$?}

\section{Алгоритм исчисления индексов}

Пусть $p$ "--- простое, $g$ "--- первообразный корень по модулю $p$, $h \in \Znx{p}$. Логарифмы
определены по модулю $n = p - 1$. Элементы $\Znx{p}$ отождествляем с их \emph{наименьшими
неотрицательными вычетами} $1, \dots, p - 1$ "--- обычными целыми числами, которые можно
раскладывать на множители.

\begin{definition}
\emph{Факторная база} "--- множество $F = \{\ell_1, \dots, \ell_t\}$ всех простых чисел
$\ell \le B$, $t = \pi(B)$. Неизвестные первого этапа "--- их логарифмы
$L_\ell = \log_g \ell \in \Zn{n}$.
\end{definition}

\begin{lemma}\label{l4:relation}
\begin{enumerate}
  \item Если $g^k \bmod p = \prod_{\ell \in F} \ell^{e_\ell}$ (равенство целых чисел), то
    \[
      \sum_{\ell \in F} e_\ell L_\ell \equiv k \pmod{p - 1}.
    \]
  \item Если $h g^{-k} \bmod p = \prod_{\ell \in F} \ell^{f_\ell}$, то
    $\log_g h \equiv k + \sum_{\ell} f_\ell L_\ell \pmod{p - 1}$.
\end{enumerate}
\end{lemma}

\begin{proof}
(1) $g^{k} \equiv \prod \ell^{e_\ell} \equiv \prod (g^{L_\ell})^{e_\ell} = g^{\sum e_\ell L_\ell}
\pmod p$, а $g^a \equiv g^b \pmod p$ равносильно $a \equiv b \pmod{\ord g}$, $\ord g = p - 1$.
(2) Аналогично $h \equiv g^{k} \prod \ell^{f_\ell} \equiv g^{k + \sum f_\ell L_\ell}$.
\end{proof}

Соотношения вида (1) называются \emph{соотношениями} (relations). Они не зависят от $h$:
первый этап выполняется один раз для данных $p$ и $g$, после чего каждый новый логарифм
вычисляется на втором, гораздо более дешёвом этапе.

\begin{algorithm}[исчисление индексов в $\Znx{p}$]\label{l4:alg}
\begin{algorithmic}[1]
\Require простое $p$, первообразный корень $g$, $h \in \Znx{p}$, граница $B$.
\Ensure $x = \log_g h \bmod (p - 1)$.
\State $F \gets \{\ell \le B : \ell \text{ простое}\}$, $t \gets \abs{F}$
\Statex \textit{Этап 1: соотношения и логарифмы факторной базы}
\Repeat
  \State выбрать $k \in [1, p - 2]$ (случайно или подряд), $a \gets g^k \bmod p$
  \If{$a$ раскладывается над $F$: $a = \prod_\ell \ell^{e_\ell}$}
    \State запомнить соотношение $\sum_\ell e_\ell L_\ell \equiv k \pmod{p - 1}$
  \EndIf
\Until{соотношений чуть больше $t$ и система однозначно разрешима}
\State решить систему по модулю $q^e$ для каждого $q^e \parallel p - 1$, собрать $L_\ell$ по КТО
\Statex \textit{Этап 2: логарифм данного элемента}
\Repeat
  \State выбрать $k$, $a \gets h g^{-k} \bmod p$
\Until{$a = \prod_\ell \ell^{f_\ell}$ раскладывается над $F$}
\State \Return $x = k + \sum_\ell f_\ell L_\ell \bmod (p - 1)$
\end{algorithmic}
\end{algorithm}

\begin{remark}
На втором этапе можно перебирать и $h g^{k}$; тогда $x \equiv \sum_\ell f_\ell L_\ell - k$.
Так сделано в задачах taskgen: $h g^{k}$ получается из $h g^{k - 1}$ одним умножением на $g$,
не нужно вычислять $g^{-1}$.
\end{remark}

\begin{remark}\label{l4:minus-one}
Часто к факторной базе добавляют $-1$ и раскладывают \emph{абсолютно наименьший} вычет $r$,
$\abs{r} < p/2$: такие числа вдвое меньше, а значит, чаще гладкие. Логарифм $-1$ известен заранее:
$g^{(p-1)/2}$ имеет порядок $2$, а единственный элемент порядка $2$ в циклической группе
$\Znx{p}$ "--- это $-1$; значит, $L_{-1} = (p - 1)/2$. Записывать $-1$ как «простое» нельзя без
этой оговорки: в соотношении знак учитывается слагаемым $L_{-1}$.
\end{remark}

\begin{remark}
В протоколах обычно $g$ порождает подгруппу простого порядка $q \mid p - 1$, а не всю $\Znx{p}$
(лекция~1). Тогда метод применяют к первообразному корню $\gamma$: если $g = \gamma^{a}$,
$h = \gamma^{b}$, то $a x \equiv b \pmod{p - 1}$, и $x$ находится по модулю $q$. Сложность
исчисления индексов определяется размером $p$, а не $q$: поэтому в $\Znx{p}$ нужны $p$ длиной
в тысячи бит, хотя $q \approx 2^{256}$ уже защищает от универсальных алгоритмов.
\end{remark}

\taskexample{ic-relations-signed}

\checkyourself{Почему в лемме~\ref{l4:relation} нужно, чтобы $g$ был первообразным корнем? Что
пойдёт не так, если $\ord g < p - 1$ и, например, $2 \notin \langle g \rangle$?}

\section{Линейная алгебра по модулю \texorpdfstring{$p - 1$}{p - 1}}

Система первого этапа "--- линейная система $E L \equiv k \pmod{n}$, $n = p - 1$, где строки
матрицы $E$ "--- векторы показателей $(e_\ell)$. Кольцо $\Zn{n}$ не является полем ($n$ чётно!),
и метод Гаусса в нём «как есть» не работает: делить можно только на обратимые элементы. Выход "---
китайская теорема об остатках (лекция~2): $\Zn{n} \cong \prod_{q^e \parallel n} \Zn{q^e}$, поэтому
систему достаточно решить по модулю каждого примарного делителя $q^e$ и собрать решения по КТО.

\begin{theorem}\label{l4:local}
Пусть $q$ "--- простое, $e \ge 1$, $A$ "--- матрица $m \times t$ над $\Zn{q^e}$, и система
$Ax = b$ имеет решение. Следующие условия равносильны:
\begin{enumerate}
  \item решение единственно по модулю $q^e$;
  \item ранг матрицы $A \bmod q$ над полем $\FF_q$ равен $t$;
  \item метод Гаусса с \emph{обратимыми} ведущими элементами (не делящимися на $q$) проходит все
    $t$ столбцов.
\end{enumerate}
\end{theorem}

\begin{proof}
(2)$\Rightarrow$(3). Элементарные преобразования строк с обратимыми множителями сохраняют ранг
по модулю $q$. Если на шаге $j$ все элементы столбца $j$ в строках $j, \dots, m$ делятся на $q$,
то по модулю $q$ первые $j$ столбцов текущей матрицы имеют ранг не больше $j - 1$ (в них
ненулевые по модулю $q$ элементы есть лишь в строках $1, \dots, j - 1$), т.\,е. $\operatorname{rk}
(A \bmod q) < t$ "--- противоречие.

(3)$\Rightarrow$(1). Метод Гаусса приводит систему к равносильной вида $x_j = c_j$,
$j = 1, \dots, t$ (плюс тождества $0 = 0$), поэтому решение единственно.

(1)$\Rightarrow$(2). Пусть $\operatorname{rk}(A \bmod q) < t$. Тогда существует вектор
$v \in \ZZ^t$, $v \not\equiv 0 \pmod q$, с $Av \equiv 0 \pmod q$. Вектор $w = q^{e-1} v$ ненулевой
по модулю $q^e$ и $Aw = q^{e-1} A v \equiv 0 \pmod{q^e}$, поэтому $x + w$ "--- второе решение.
\end{proof}

Таким образом, «вырожденность» системы надо проверять по модулю \emph{каждого} простого
делителя $q$ числа $p - 1$. Для квадратной системы это условие $q \nmid \det E$. Например,
соотношение $5 L_3 \equiv 15 \pmod{700}$ определяет $L_3$ по модулю $4$ и $7$, но по модулю $25$
даёт лишь $L_3 \equiv 3 \pmod 5$: коэффициент $5$ не обратим, и нужны другие соотношения.

\begin{proposition}\label{l4:random-rank}
Вероятность того, что $t$ независимых равномерно распределённых векторов из $\FF_q^t$ линейно
независимы, равна $\prod_{i=1}^{t}(1 - q^{-i})$; при $q = 2$ она больше $0{,}288$.
Вероятность того, что $t + s$ таких векторов порождают $\FF_q^t$, не меньше $1 - q^{-s}/(q - 1)$.
\end{proposition}

\begin{proof}
Первые $i - 1$ независимых векторов порождают подпространство из $q^{i-1}$ элементов; $i$-й
вектор не лежит в нём с вероятностью $1 - q^{i-1}/q^t$. Перемножая по $i = 1, \dots, t$, получаем
$\prod_{j=1}^{t}(1 - q^{-j})$; при $q = 2$ это произведение убывает по $t$ к
$\prod_{j \ge 1}(1 - 2^{-j}) = 0{,}2887\ldots$. Для второго утверждения: $t + s$ векторов
не порождают $\FF_q^t$ тогда и только тогда, когда все они лежат в некоторой гиперплоскости
$\{x : a \cdot x = 0\}$, $a \ne 0$. Гиперплоскостей $(q^t - 1)/(q - 1)$, и каждая содержит все
векторы с вероятностью $q^{-(t+s)}$; по неравенству для вероятности объединения вероятность
неудачи не больше $\frac{q^t - 1}{q - 1} q^{-t-s} < \frac{q^{-s}}{q - 1}$.
\end{proof}

Векторы показателей, конечно, не равномерны (большие показатели при $2$ встречаются чаще, чем
при $7$, а многие соотношения содержат лишь несколько малых простых), но вывод верен и на
практике: соотношений нужно собрать $t + s$ с небольшим запасом $s$ (порядка $10$--$20$).
По модулю $2$ вырожденность встречается чаще всего "--- и именно её легко обойти.

\begin{proposition}\label{l4:parity}
Пусть $g$ "--- первообразный корень по модулю нечётного простого $p$, $\ell \in \Znx{p}$. Тогда
$\log_g \ell$ чётен тогда и только тогда, когда символ Лежандра $\legendre{\ell}{p} = 1$.
Вообще, для $q^e \parallel p - 1$ значение $\log_g \ell \bmod q^e$ "--- это дискретный логарифм
$\ell^{(p-1)/q^e}$ по основанию $g^{(p-1)/q^e}$ в группе порядка $q^e$.
\end{proposition}

\begin{proof}
По критерию Эйлера $\legendre{\ell}{p} \equiv \ell^{(p-1)/2} = g^{L_\ell (p-1)/2} =
(g^{(p-1)/2})^{L_\ell} = (-1)^{L_\ell}$, так как $g^{(p-1)/2} = -1$ (замечание~\ref{l4:minus-one}).
Второе утверждение "--- редукция Полига"--~Хеллмана (лекция~2): $(g^{(p-1)/q^e})^{L_\ell} =
\ell^{(p-1)/q^e}$, а $g^{(p-1)/q^e}$ имеет порядок $q^e$.
\end{proof}

Поэтому на практике по модулю малых простых делителей $p - 1$ логарифмы факторной базы находят
напрямую (символами Лежандра, алгоритмом Полига"--~Хеллмана), а линейную систему решают только
по модулю большого простого делителя $q$ числа $p - 1$ "--- над \emph{полем} $\FF_q$. Для
«безопасного» простого $p = 2q + 1$ (лекция~1) это вся линейная алгебра.

\taskexample{ic-system-singular}

\checkyourself{Соотношения $g^{4} \equiv 2^2 \cdot 3^2$, $g^{6} \equiv 2^4$, $g^{10} \equiv 3^2
\cdot 5^2 \pmod p$ "--- что они дают по модулю $2$? Что можно сказать о $\legendre{2}{p}$, если
среди соотношений есть $g^{7} \equiv 2^{3} \cdot 3^{2}$?}

\section{Пример}

Пусть $p = 701$, $p - 1 = 700 = 2^2 \cdot 5^2 \cdot 7$, $g = 11$ "--- первообразный корень:
$11^{350} \equiv 700$, $11^{140} \equiv 89$, $11^{100} \equiv 550$, т.\,е. $11^{700/q} \not\equiv 1$
для $q = 2, 5, 7$. Возьмём $B = 7$, $F = \{2, 3, 5, 7\}$, $t = 4$.
(Рецепт следующего раздела $B \approx L(p)^{1/\sqrt2}$ даёт $B \approx 12$; мы берём меньшую
границу, чтобы система была короче.)

\emph{Соотношения.} Среди $11^k \bmod 701$, $k = 1, \dots, 29$, гладкими оказываются семь:
\[
\begin{array}{r@{\;\equiv\;}l@{\qquad\Rightarrow\qquad}l}
  g^{3} & 630 = 2 \cdot 3^2 \cdot 5 \cdot 7 & L_2 + 2L_3 + L_5 + L_7 \equiv 3,\\
  g^{7} & 72 = 2^3 \cdot 3^2 & 3L_2 + 2L_3 \equiv 7,\\
  g^{9} & 300 = 2^2 \cdot 3 \cdot 5^2 & 2L_2 + L_3 + 2L_5 \equiv 9,\\
  g^{15} & 243 = 3^5 & 5L_3 \equiv 15,\\
  g^{22} & 672 = 2^5 \cdot 3 \cdot 7 & 5L_2 + L_3 + L_7 \equiv 22,\\
  g^{28} & 320 = 2^6 \cdot 5 & 6L_2 + L_5 \equiv 28,\\
  g^{29} & 15 = 3 \cdot 5 & L_3 + L_5 \equiv 29
\end{array}
\pmod{700}.
\]
(Доля $7/29 \approx 0{,}24$ сравнима с долей $7$-гладких чисел до $700$: $\psi(700, 7)/700 \approx
0{,}17$.)

\emph{По модулю $4$.} Из $(15)$: $5 \equiv 1$, так что $L_3 \equiv 15 \equiv 3$. Из $(29)$:
$L_5 \equiv 29 - 3 \equiv 2$. Соотношение $(28)$ даёт $6L_2 \equiv 28 - 2$, т.\,е. $2L_2 \equiv 2
\pmod 4$ "--- лишь нечётность $L_2$ (коэффициент $2$ не обратим). Берём $(7)$: $3L_2 \equiv 7 -
2 \cdot 3 = 1$, $L_2 \equiv 3 \cdot 1 = 3$ ($3^{-1} \equiv 3 \pmod 4$). Из $(22)$:
$L_7 \equiv 22 - 5 \cdot 3 - 3 = 4 \equiv 0$.

\emph{По модулю $7$.} $(15)$: $5L_3 \equiv 1$, $L_3 \equiv 3$ ($5^{-1} \equiv 3$). $(29)$:
$L_5 \equiv 29 - 3 \equiv 5$. $(28)$: $6L_2 \equiv 28 - 5 \equiv 2$, $L_2 \equiv 2 \cdot 6 \equiv 5$
($6^{-1} \equiv 6$). $(22)$: $L_7 \equiv 22 - 25 - 3 \equiv 1$.

\emph{По модулю $25$.} Здесь $(15)$ даёт только $L_3 \equiv 3 \pmod 5$. Исключим $L_5 \equiv 4 -
L_3$ с помощью $(29)$; тогда $(28)$: $6L_2 - L_3 \equiv -1$, т.\,е. $L_3 \equiv 6L_2 + 1$, и $(7)$:
$3L_2 + 2(6L_2 + 1) \equiv 7$, т.\,е. $15 L_2 \equiv 5 \pmod{25}$ "--- снова необратимый
коэффициент, $L_2 \equiv 2 \pmod 5$. Соотношения $(7), (15), (28), (29)$ вырождены по модулю $5$
(по модулю $5$ соотношение $(15)$ тривиально, а матрица соотношений $(7), (28), (29)$
относительно $L_2, L_3, L_5$ имеет определитель $-15$, кратный $5$), и нужно ещё одно. Подставляя в $(9)$:
$2L_2 + L_3 + 2(4 - L_3) = 2L_2 - L_3 + 8 \equiv 9$, откуда $2L_2 - 6L_2 - 1 \equiv 1$,
$4L_2 \equiv -2 \equiv 23$ и $L_2 \equiv 23 \cdot 19 \equiv 12$ ($4^{-1} \equiv 19$). Далее
$L_3 \equiv 73 \equiv 23$, $L_5 \equiv 4 - 23 \equiv 6$, а из $(22)$: $L_7 \equiv 22 - 60 - 23 \equiv 14$.

\emph{КТО} по модулям $4, 25, 7$:
\[
\begin{array}{c|ccc|r}
  & \bmod 4 & \bmod 25 & \bmod 7 & L_\ell \\ \hline
  \ell = 2 & 3 & 12 & 5 & 187\\
  \ell = 3 & 3 & 23 & 3 & 423\\
  \ell = 5 & 2 & 6 & 5 & 306\\
  \ell = 7 & 0 & 14 & 1 & 64
\end{array}
\]
Проверка чётности по утверждению~\ref{l4:parity}: $701 \equiv 5 \pmod 8$, поэтому
$\legendre{2}{701} = -1$, и $L_2$ нечётен; $\legendre{3}{701} = \legendre{701}{3} = \legendre{2}{3}
= -1$ "--- $L_3$ нечётен; $\legendre{5}{701} = \legendre{1}{5} = 1$, $\legendre{7}{701} =
\legendre{1}{7} = 1$ "--- $L_5$, $L_7$ чётны. Неиспользованное соотношение $(3)$ "--- проверочное:
$187 + 846 + 306 + 64 = 1403 \equiv 3 \pmod{700}$.

\emph{Второй этап}: найдём $\log_{11} 123$. Здесь $g^{-1} \equiv 255$, и
\[
\begin{gathered}
  123 = 3 \cdot 41,\quad 123 g^{-1} \equiv 521,\quad 123 g^{-2} \equiv 366 = 2 \cdot 3 \cdot 61,\\
  123 g^{-3} \equiv 97,\quad 123 g^{-4} \equiv 200 = 2^3 \cdot 5^2.
\end{gathered}
\]
По лемме~\ref{l4:relation} $\log_{11} 123 \equiv 4 + 3 L_2 + 2 L_5 = 4 + 561 + 612 = 1177 \equiv 477
\pmod{700}$; действительно, $11^{477} \equiv 123 \pmod{701}$.

\taskexample{ic-full}

\taskexample{ic-phase2-signed}

\checkyourself{В примере для $p = 701$ найдите $\log_{11} 2026$ (т.\,е. логарифм $2026 \bmod 701
= 624$), пользуясь найденными $L_\ell$. Сколько попыток пришлось сделать?}

\section{Оценка сложности}

Анализ метода опирается на эвристическое предположение: \emph{вычеты $g^k \bmod p$ (и $h g^{-k}
\bmod p$) при случайном $k$ ведут себя по отношению к гладкости как случайные числа из
$[1, p)$}, т.\,е. гладкие с вероятностью $\psi(p, B)/p$. (Для строгого варианта алгоритма
К.~Померанс доказал ту же итоговую оценку $L_p[1/2, \sqrt 2]$ без эвристик.)

Пусть $B = L_p[1/2, c]$, $t = \pi(B) = B^{1 + o(1)} = L_p[1/2, c]$ (лемма~\ref{l4:Lprops}:
множитель $1/\ln B$ поглощается). Подсчитаем стоимость этапов.

\begin{enumerate}
  \item \emph{Число попыток.} Нужно $t + \bigO(1)$ соотношений
    (утверждение~\ref{l4:random-rank}), а каждое случайное $g^k$ гладкое с вероятностью
    $L_p[1/2, -\frac{1}{2c}]$ (лемма~\ref{l4:uu}). Ожидаемое число попыток
    \[
      N = t \cdot \frac{p}{\psi(p, B)} = L_p\Bigl[\tfrac12, c + \tfrac{1}{2c}\Bigr].
    \]
  \item \emph{Стоимость попытки}: одно умножение по модулю $p$ (полиномиально по $\ln p$) и
    проверка гладкости. Пробное деление стоит $\pi(B) = L_p[1/2, c]$ делений; существуют и
    методы проверки гладкости стоимости $L_p[1/2, o(1)]$, т.\,е. $p^{o(1)}$ в $L$-смысле, "---
    например, метод эллиптических кривых Ленстры (лекция~10).
  \item \emph{Линейная алгебра}: система $\approx t \times t$. Метод Гаусса "--- $\bigO(t^3)
    = L_p[1/2, 3c]$ операций. Но система \emph{разрежена}: в строке не больше $\log_2 p$
    ненулевых элементов, и методы Ланцоша и Видемана решают её за $t^{2 + o(1)} = L_p[1/2, 2c]$.
  \item \emph{Второй этап}: $p/\psi(p, B) = L_p[1/2, \frac{1}{2c}]$ попыток.
\end{enumerate}

\begin{theorem}[эвристическая оценка]\label{l4:complexity}
При выборе $B = L_p[1/2, 1/\sqrt2]$, т.\,е. $B \approx L(p)^{1/\sqrt 2}$, быстрой проверке
гладкости и разреженной линейной алгебре первый этап метода исчисления индексов выполняется за
$L_p[1/2, \sqrt 2]$ операций, а второй "--- за $L_p[1/2, 1/\sqrt2]$; это минимум по $c$.
С пробным делением оптимальны $c = 1/2$ и сложность $L_p[1/2, 2]$ (второй этап "---
$L_p[1/2, 3/2]$).
\end{theorem}

\begin{proof}
По лемме~\ref{l4:Lprops} показатели стоимостей складываются при умножении, а при сложении берётся
максимум. С быстрой проверкой гладкости и разреженной алгеброй показатель первого этапа равен
\[
  f(c) = \max\Bigl(c + \frac{1}{2c},\ 2c\Bigr).
\]
Функция $c + \frac{1}{2c}$ имеет минимум $\sqrt 2$ при $c = 1/\sqrt 2$ (производная $1 -
\frac{1}{2c^2}$ меняет знак с минуса на плюс), а $2c$ возрастает. При $c \le 1/\sqrt2$ имеем
$f(c) \ge c + \frac{1}{2c} \ge \sqrt2$, при $c \ge 1/\sqrt2$ "--- $f(c) \ge 2c \ge \sqrt 2$,
и $f(1/\sqrt2) = \max(\sqrt2, \sqrt2) = \sqrt 2$. Второй этап стоит $L_p[1/2, \frac{1}{2c}] =
L_p[1/2, \frac{1}{\sqrt2}]$.

С пробным делением каждая попытка дороже в $L_p[1/2, c]$ раз:
$f(c) = \max\bigl(2c + \frac{1}{2c}, 2c\bigr) = 2c + \frac{1}{2c}$, минимум $2$ при $c = 1/2$
(производная $2 - \frac{1}{2c^2}$); второй этап стоит
$L_p[1/2, c + \frac{1}{2c}] = L_p[1/2, \frac32]$.
\end{proof}

\begin{remark}
Учебники (в частности, \cite[разд.~3.8]{hps}) часто учитывают лишь число попыток $N =
L_p[1/2, c + \frac{1}{2c}]$ и получают тот же рецепт $B = L(p)^{1/\sqrt2}$ и ту же оценку
$L(p)^{\sqrt 2}$. С плотным методом Гаусса показатель равен $\max(c + \frac{1}{2c}, 3c)$; его
минимум $\frac32$ достигается при $c = \frac12$ "--- линейная алгебра становится узким местом,
поэтому на практике так важны разреженные методы.
\end{remark}

Без $o(1)$ эти выражения дают лишь грубое представление о масштабе, но сравнение показательно
(логарифмы по основанию $2$ числа операций):
\[
\begin{array}{r|r|r|r}
  \log_2 p & \sqrt{p} & L_p[1/2, \sqrt 2] & L_p[1/3, 1{,}923] \\ \hline
  512 & 256 & 93 & 64 \\
  1024 & 512 & 139 & 87 \\
  2048 & 1024 & 207 & 117 \\
  3072 & 1536 & 261 & 139
\end{array}
\]
Последний столбец "--- решето числового поля (см. ниже). Именно субэкспоненциальные алгоритмы
определяют рекомендуемые длины $p$: $2048$ бит и больше.

\checkyourself{Почему при оптимальном выборе $B$ стоимости поиска соотношений и линейной
алгебры совпадают? Как изменится оптимальное $c$, если бы линейная алгебра стоила $t^{4}$?}

\section{Современные варианты и пределы применимости}

\begin{itemize}
  \item \emph{Решето числового поля} (NFS, Д.~Гордон~\cite{gordon-nfs-dlog}, 1993) заменяет
    кольцо $\ZZ$ кольцами целых алгебраических чисел, где «гладкими» должны быть числа
    существенно меньшего размера, и достигает сложности $L_p[1/3, (64/9)^{1/3}]$,
    $(64/9)^{1/3} \approx 1{,}923$. Структура та же: сбор соотношений, линейная алгебра по модулю
    большого делителя $p - 1$, «спуск» для отдельного логарифма. В 2019 г. так был
    вычислен дискретный логарифм по модулю $240$-значного ($795$-битного) простого~\cite{boudot-240}.
  \item \emph{Предвычисление} зависит только от $p$. Атака Logjam~\cite{logjam} (2015) показала,
    что для 512-битных простых, общих для множества серверов, первый этап выполняется один раз,
    после чего каждый логарифм находится примерно за минуту. Отсюда рекомендация: $p \ge 2048$ бит и
    не использовать широко распространённые малые модули.
  \item В полях $\FF_{2^n}$, $\FF_{3^n}$ (малая характеристика) варианты исчисления индексов
    после 2013 г. стали \emph{квазиполиномиальными}; такие поля в криптографии больше
    не применяют.
  \item \emph{Эллиптические кривые} (лекции 8--9). Для группы точек $\EC(\FF_p)$ общего вида
    аналог метода не известен: у точек нет «размера», согласованного с групповой операцией,
    и нет разложения на «малые простые точки». Попытки поднять точки на кривую над $\QQ$
    наталкиваются на быстрый рост высот, а гладкие точки не встречаются достаточно часто
    (В.~Миллер~\cite{miller-ecc}). Поэтому для ECDLP лучшие известные алгоритмы "---
    универсальные, $\bigO(\sqrt q)$, и 256-битной кривой хватает там, где в $\Znx{p}$ нужен
    3072-битный модуль. Исключения (сверхсингулярные и аномальные кривые) сводят ECDLP к DLP
    в конечном поле или решают его иначе; они обсуждаются в лекции~9.
\end{itemize}

Общая схема исчисления индексов в группе $G$ требует: (1) небольшого множества $F \subset G$
«малых» элементов; (2) заметной доли элементов $G$, раскладывающихся в произведение элементов $F$;
(3) эффективного алгоритма разложения. В $\Znx{p}$ все три условия даёт арифметика $\ZZ$;
для эллиптических кривых общего вида не выполнено уже (1)--(2).

\checkyourself{Какие из условий (1)--(3) нарушаются для группы $\Zn{n}$ по сложению
(DLP в ней "--- это решение $xg \equiv h \pmod n$)? Почему там исчисление индексов и не нужно?}

\checkyourselfsection

\begin{furtherreading}
  \readitem[разд.~3.7--3.8]{hps}{гладкие числа, функция $\psi(x, B)$, $L$-нотация и метод
    исчисления индексов с подробным примером}
  \readitem[разд.~1.4.5, 6.4]{crandall-pomerance}{гладкие числа; исчисление индексов для
    дискретного логарифма}
  \readitem[гл.~15]{galbraith}{субэкспоненциальные алгоритмы: гладкие числа, исчисление индексов в
    конечных полях, решето числового поля}
  \readitem[гл.~15]{shoup}{строгий анализ простого варианта исчисления индексов}
  \readitem[гл.~3]{vasilenko}{логарифмирование в конечных полях (на русском языке)}
  \readitem{odlyzko-dlog}{обзор истории и состояния алгоритмов дискретного логарифмирования}
  \readitem{adleman-index-calculus}{первый субэкспоненциальный алгоритм для DLP}
  \readitem{canfield-erdos-pomerance}{оценка числа гладких чисел}
  \readitem{logjam}{практическая атака на Диффи"--~Хеллмана с 512-битными простыми}
  \readitem[конспект «Index calculus, smooth numbers, and factoring integers»]{mit-18783}{лекция
    курса MIT с доказательством оценки для простого варианта алгоритма}
\end{furtherreading}
\end{lecture}
